\NormalFont\ShortTitle{Efeseni}

{\MT Scrisoarea lui Pavel către Efeseni


\par }\ChapOne{1}{\PP \VerseOne{1}Pavel, apostol al lui Isus Hristos, prin voia lui Dumnezeu, către sfinții care sunt în Efes și către cei credincioși în Isus Hristos:

\VS{2}Har vouă și pace de la Dumnezeu Tatăl nostru și de la Domnul Isus Hristos.


\par }{\PP \VS{3}Binecuvântat să fie Dumnezeul și Tatăl Domnului nostru Isus Hristos, care ne-a binecuvântat cu orice binecuvântare duhovnicească în locurile cerești, în Hristos,

\VS{4}după cum ne-a ales în El înainte de întemeierea lumii, ca să fim sfinți și fără cusur înaintea Lui, în dragoste,

\VS{5}predestinându-ne pentru adopție, ca copii, prin Isus Hristos, pentru El însuși, după buna plăcere a dorinței Sale,

\VS{6}spre lauda gloriei harului Său, prin care ne-a dat de bunăvoie favoarea în Cel iubit.

\VS{7}În El avem răscumpărarea noastră prin sângele Său, iertarea greșelilor noastre, potrivit cu bogăția haruluiSău

\VS{8}pe care l-a făcut să abunde față de noi în toată înțelepciunea și prudența,

\VS{9}făcându-ne cunoscut misterul voinței Sale, potrivit bunei Sale plăceri pe care a hotărât-o în El

\VS{10}pentru o administrare a plinătății vremurilor, pentru a rezuma toate lucrurile în Hristos, cele din ceruri și cele de pe pământ, în El.

\VS{11}Ni s-a atribuit și nouă o moștenire în El, fiind rânduiți mai dinainte, potrivit cu planul Celui care face toate lucrurile după planul voii Sale,

\VS{12}pentru ca noi, care mai înainte am sperat în Hristos, să fim spre lauda gloriei Lui.

\VS{13}În el și voi, după ce ați auzit cuvântul adevărului, Vestea cea bună a mântuirii voastre — în care, crezând și voi, ați fost pecetluiți cu Duhul Sfânt făgăduit,

\VS{14}care este un gaj al moștenirii noastre, pentru răscumpărarea proprietății lui Dumnezeu, spre lauda gloriei sale.


\par }{\PP \VS{15}De aceea și eu, care am auzit despre credința în Domnul Isus, care este printre voi, și despre dragostea pe care o aveți față de toți sfinții,

\VS{16}nu încetez să mulțumesc pentru voi, pomenindu-vă în rugăciunile mele,

\VS{17}pentru ca Dumnezeul Domnului nostru Isus Hristos, Tatăl slavei, să vă dea un duh de înțelepciune și de descoperire în cunoașterea Lui,

\VS{18}ca să vă lumineze ochii inimii voastre, ca să știți care este nădejdea chemării Lui, și care sunt bogățiile gloriei moștenirii Sale în sfinți,

\VS{19}și care este nespus de mare puterea Sa față de noi, cei care credem, potrivit cu lucrarea tăriei puterii Sale

\VS{20}pe care a lucrat-o în Hristos când L-a înviat din morți și L-a făcut să șadă la dreapta Sa în locurile cerești,

\VS{21}cu mult mai presus de orice stăpânire, autoritate, putere, domnie și orice nume care se numește, nu numai în acest veac, ci și în cel ce va să vină.

\VS{22}I-a supus toate lucrurile sub picioarele Lui și L-a dat să fie cap peste toate lucrurile pentru Adunare,

\VS{23}care este trupul Lui, plinătatea Celui care umple totul în toți.



\par }\Chap{2}{\PP \VerseOne{1}Voi ați fost înviați, când erați morți în fărădelegi și păcate,

\VS{2}în care umblați odinioară, după cursul acestei lumi, după prințul puterii aerului, duhul care lucrează acum în copiii neascultătorilor.

\VS{3}Și noi toți am trăit odinioară printre ei în poftele cărnii noastre, împlinind poftele cărnii și ale minții, și eram din fire copii ai mâniei, ca și ceilalți.

\VS{4}Dar Dumnezeu, fiind bogat în îndurare, pentru marea Sa dragoste cu care ne-a iubit,

\VS{5}chiar și atunci când eram morți din cauza greșelilor noastre, ne-a făcut vii împreună cu Hristos — prin har ați fost mântuiți —

\VS{6}și ne-a înviat împreună cu El și ne-a făcut să ședem cu El în locurile cerești în Hristos Isus,

\VS{7}pentru ca în veacurile viitoare să arate bogăția nespus de mare a harului Său în bunătatea Sa față de noi în Hristos Isus;

\VS{8}căci prin har ați fost mântuiți prin credință, și aceasta nu de la voi înșivă; este un dar al lui Dumnezeu,

\VS{9}nu prin fapte, ca nimeni să nu se laude.

\VS{10}Căci noi suntem lucrarea Lui, creați în Hristos Isus pentru fapte bune, pe care Dumnezeu le-a pregătit mai dinainte ca noi să umblăm în ele.


\par }{\PP \VS{11}Aduceți-vă aminte că odinioară, voi, neamurile în trup, care sunteți numiți “netăiați împrejur” prin ceea ce se numește “tăierea împrejur” (în trup, făcută de mâini),

\VS{12}erați în vremea aceea despărțiți de Hristos, înstrăinați de comunitatea lui Israel și străini de legămintele făgăduinței, fără nădejde și fără Dumnezeu în lume.

\VS{13}Dar acum, în Hristos Isus, voi, care altădată erați departe, ați fost apropiați prin sângele lui Hristos.

\VS{14}Căci el este pacea noastră, care i-a făcut pe amândoi una și a dărâmat zidul de mijloc al despărțirii,

\VS{15}după ce a desființat în trupul său ostilitatea, legea poruncilor cuprinse în rânduieli, ca să creeze în sine un singur om nou din cei doi, făcând pace,

\VS{16}și să-i împace pe amândoi într-un singur trup cu Dumnezeu prin cruce, după ce a ucis prin ea ostilitatea.

\VS{17}El a venit și a propovăduit pacea atât vouă, care erați departe, cât și celor apropiați.

\VS{18}Căci prin el amândoi avem acces la Tatăl într-un singur Duh.

\VS{19}Așadar, voi nu mai sunteți străini și străini, ci sunteți concetățeni ai sfinților și ai casei lui Dumnezeu,

\VS{20}fiind zidiți pe temelia apostolilor și a profeților, piatra unghiulară principală fiind Cristos Isus însuși;

\VS{21}în care toată clădirea, bine închegată, se dezvoltă într-un templu sfânt în Domnul;

\VS{22}în care și voi sunteți zidiți împreună pentru a fi locuința lui Dumnezeu în Duhul Sfânt.



\par }\Chap{3}{\PP \VerseOne{1}De aceea, eu, Pavel, sunt prizonierul lui Hristos Isus pentru voi, neamurilor,

\VS{2}dacă ați auzit despre administrarea harului lui Dumnezeu, care mi-a fost dat față de voi,

\VS{3}și despre faptul că, prin revelație, mi s-a făcut cunoscut misterul, așa cum am scris mai înainte în puține cuvinte,

\VS{4}și prin care, când citiți, puteți pricepe înțelegerea mea în misterul lui Hristos,

\VS{5}care în alte generații nu a fost făcut cunoscut fiilor oamenilor, așa cum a fost descoperit acum sfinților Săi apostoli și profeți în Duhul Sfânt,

\VS{6}că neamurile sunt părtașe la moștenire și la trup și părtașe la făgăduința Lui în Hristos Isus prin Buna Vestire,

\VS{7}de care am fost făcut slujitor, potrivit darului harului lui Dumnezeu care mi-a fost dat, după lucrarea puterii Lui.

\VS{8}Mie, cel mai mic dintre toți sfinții, mi-a fost dat acest har, ca să propovăduiesc neamurilor bogățiile nepătrunse ale lui Hristos

\VS{9}și să fac pe toți oamenii să vadă care este administrarea misterului care de veacuri a fost ascuns în Dumnezeu, care a creat toate lucrurile prin Isus Hristos,

\VS{10}pentru ca acum, prin adunare, înțelepciunea multiplă a lui Dumnezeu să fie făcută cunoscută principatelor și puterilor din locurile cerești,

\VS{11}conform scopului veșnic pe care l-a împlinit în Hristos Isus, Domnul nostru.

\VS{12}În El avem îndrăzneală și acces cu încredere, prin credința noastră în El.

\VS{13}De aceea vă rog să nu vă pierdeți inima din cauza necazurilor mele pentru voi, care sunt gloria voastră.


\par }{\PP \VS{14}De aceea, îmi plec genunchii înaintea Tatălui Domnului nostru Isus Hristos,

\VS{15}de la care se numește orice familie în ceruri și pe pământ,

\VS{16}ca să vă dea, după bogăția slavei Lui, să fiți întăriți cu putere, prin Duhul Său, în lăuntrul vostru,

\VS{17}pentru ca Hristos să locuiască în inimile voastre prin credință, pentru ca voi, fiind înrădăcinați și întemeiați în dragoste,

\VS{18}să fiți întăriți ca să înțelegeți împreună cu toți sfinții care este lățimea, lungimea, înălțimea și adâncimea,

\VS{19}și să cunoașteți dragostea lui Hristos, care întrece orice cunoștință, ca să fiți plini de toată plinătatea lui Dumnezeu.


\par }{\PP \VS{20}Iar Celui ce poate să facă mai mult decât cerem și decât gândim noi, după puterea care lucrează în noi,

\VS{21}Lui să fie slava în adunare și în Hristos Isus, în toate generațiile, în vecii vecilor. Amin.



\par }\Chap{4}{\PP \VerseOne{1}De aceea, eu, prizonierul în Domnul, vă rog să vă purtați în mod demn de chemarea cu care ați fost chemați,

\VS{2}cu toată smerenia și umilința, cu răbdare, suportându-vă unii pe alții în dragoste,

\VS{3}urmărind să păstrați unitatea Duhului în legătura păcii.

\VS{4}Există un singur trup și un singur Duh, așa cum și voi ați fost chemați într-o singură nădejde a chemării voastre,

\VS{5}un singur Domn, o singură credință, un singur botez,

\VS{6}un singur Dumnezeu și Tatăl tuturor, care este peste toate și prin toate și în noi toți.

\VS{7}Dar fiecăruia dintre noi i-a fost dat harul, după măsura darului lui Hristos.

\VS{8}De aceea spune el,

\par }{\Q “Când s-a înălțat la înălțime,

\par }{\QB a dus captivi pe cei din captivitate,

\par }{\QB și dădea daruri oamenilor.”


\par }{\MM \VS{9}Și când zice: “S-a înălțat”, ce înseamnă aceasta, dacă nu că mai întâi S-a coborât în părțile de jos ale pământului?

\VS{10}Cel care a coborât este cel care s-a înălțat și el mult deasupra tuturor cerurilor, ca să umple toate lucrurile.


\par }{\PP \VS{11}Pe unii i-a dat ca apostoli, pe alții ca profeți, pe alții ca evangheliști, pe alții ca păstori și învățători;

\VS{12}pentru desăvârșirea sfinților, pentru lucrarea de slujire, pentru zidirea trupului lui Hristos,

\VS{13}până când vom ajunge cu toții la unitatea credinței și a cunoașterii Fiului lui Dumnezeu, la un om matur, la măsura statura plinătății lui Hristos,

\VS{14}ca să nu mai fim copii, aruncați încoace și încolo și purtați de orice vânt de învățătură, prin viclenia oamenilor, cu viclenie, după vicleșugurile rătăcirii;

\VS{15}ci, vorbind adevărul în dragoste, să creștem în toate lucrurile în cel care este capul, Hristos,

\VS{16}din care tot trupul, fiind potrivit și împletit prin ceea ce furnizează fiecare încheietură, potrivit cu lucrarea pe măsură a fiecărei părți în parte, face ca trupul să crească spre zidirea lui însuși în dragoste.


\par }{\PP \VS{17}De aceea spun aceasta și mărturisesc în Domnul, ca să nu mai umblați cum umblă și celelalte neamuri, în deșertăciunea minții lor,

\VS{18}fiind întunecați în înțelegerea lor, înstrăinați de viața lui Dumnezeu, din pricina ignoranței care este în ei, din pricina împietririi inimilor lor.

\VS{19}Ei, devenind insensibili, s-au dat pe ei înșiși poftei, pentru a lucra cu lăcomie orice necurăție.

\VS{20}Dar voi nu așa l-ați învățat pe Cristos,

\VS{21}dacă, într-adevăr, l-ați auzit și ați fost învățați în el, așa cum este adevărul în Isus:

\VS{22}să vă lepădați, în ceea ce privește modul vostru de viață de mai înainte, de omul cel vechi, care se strică după poftele înșelăciunii,

\VS{23}și să vă înnoiți în duhul minții voastre,

\VS{24}și să vă îmbrăcați cu omul cel nou, care, după asemănarea lui Dumnezeu, a fost creat în neprihănire și în sfințenia adevărului.


\par }{\PP \VS{25}De aceea, lăsând deoparte minciuna, să vorbească fiecare cu adevărul cu aproapele său, pentru că suntem mădulare unul altuia.

\VS{26}“Mâniați-vă și nu păcătuiți.” Nu lăsa să apună soarele peste mânia ta,

\VS{27}și nu da loc diavolului.

\VS{28}Cel care a furat să nu mai fure, ci mai degrabă să muncească, producând cu mâinile sale ceva bun, ca să aibă ce să dea celui care are nevoie.

\VS{29}Să nu iasă din gura voastră nici un cuvânt corupt, ci numai ceea ce este bun pentru zidirea altora, după nevoie, ca să dea har celor ce ascultă.

\VS{30}Nu întristați Duhul Sfânt al lui Dumnezeu, în care ați fost pecetluiți pentru ziua răscumpărării.

\VS{31}Să se îndepărteze de la voi orice amărăciune, mânie, furie, mânie, scandal și calomnie, împreună cu orice răutate.

\VS{32}Și fiți buni unii cu alții, blânzi la inimă, iertându-vă unii pe alții, după cum și Dumnezeu v-a iertat pe voi în Hristos.



\par }\Chap{5}{\PP \VerseOne{1}Fiți deci imitatori ai lui Dumnezeu, ca niște copii iubiți.

\VS{2}Umblați în dragoste, după cum și Hristos ne-a iubit pe noi și S-a dat pe Sine însuși pentru noi, ca jertfă și ca sacrificiu pentru Dumnezeu, ca o mireasmă de bun miros.


\par }{\PP \VS{3}Dar nici măcar să nu se pomenească între voi, cum se cuvine sfinților, de desfrânare sexuală și de orice necurăție sau lăcomie,

\VS{4}nici de spurcăciune, nici de vorbărie prostească, nici de glume, care nu se cuvin, ci de mulțumiri.


\par }{\PP \VS{5}Să știi că nici un desfrânat, nici un necurat, nici un lacom, nici un idolatru nu are moștenire în Împărăția lui Hristos și a lui Dumnezeu.


\par }{\PP \VS{6}Nimeni să nu vă amăgească cu vorbe deșarte, căci din pricina acestor lucruri vine mânia lui Dumnezeu peste copiii neascultătorilor.

\VS{7}De aceea, nu fiți părtași la ei.

\VS{8}Căci odinioară erați întuneric, dar acum sunteți lumină în Domnul. Umblați ca niște copii ai luminii,

\VS{9}căci roada Duhului este în toată bunătatea, dreptatea și adevărul,

\VS{10}dovedind ceea ce este plăcut Domnului.

\VS{11}Nu aveți nici o părtășie cu faptele fără roade ale întunericului, ci mai degrabă chiar le mustrați.

\VS{12}Căci este o rușine chiar să vorbești despre lucrurile pe care le fac ei în ascuns.

\VS{13}Dar toate lucrurile, când sunt mustrate, sunt descoperite de lumină, căci tot ce descoperă este lumină.

\VS{14}De aceea zice: “Deșteaptă-te, tu, care dormi, și scoală-te din morți, și Hristos va străluci peste tine.”


\par }{\PP \VS{15}De aceea, păziți-vă bine cum umblați, nu ca niște neînțelepți, ci ca niște înțelepți,

\VS{16}răscumpărând timpul, căci zilele sunt rele.

\VS{17}De aceea, nu fiți nebuni, ci înțelegeți care este voia Domnului.

\VS{18}Nu vă îmbătați cu vin, în care este risipire, ci umpleți-vă de Duhul Sfânt,

\VS{19}vorbindu-vă unul altuia în psalmi, imnuri și cântări spirituale, cântând și făcând în inima voastră cântece de mulțumire Domnului,

\VS{20}mulțumind întotdeauna lui Dumnezeu, Tatăl, în numele Domnului nostru Isus Hristos, cu privire la toate lucrurile,

\VS{21}supunându-vă unul altuia în frica lui Hristos.


\par }{\PP \VS{22}Neveste, supuneți-vă bărbaților voștri, ca Domnului.

\VS{23}Căci soțul este capul soției, după cum și Hristos este capul adunării, fiind El însuși mântuitorul trupului.

\VS{24}Dar, după cum Adunarea este supusă lui Hristos, tot așa și soțiile să fie supuse în totul propriilor lor soți.


\par }{\PP \VS{25}Bărbaților, iubiți-vă nevestele, după cum și Hristos a iubit Adunarea și S-a dat pe Sine însuși pentru ea,

\VS{26}ca să o sfințească, după ce a curățat-o prin spălarea cu apă și cu cuvântul,

\VS{27}ca să Se înfățișeze pe Sine însuși în chip glorios, fără pată, fără riduri sau altceva de acest fel, ci ca să fie sfântă și fără cusur.

\VS{28}Tot așa și soții trebuie să iubească pe soțiile lor ca pe propriile lor trupuri. Cel care își iubește propria soție se iubește pe sine însuși.

\VS{29}Căci nimeni nu și-a urât vreodată propriul trup, ci îl hrănește și îl îngrijește, așa cum face și Domnul cu adunarea,

\VS{30}pentru că suntem mădulare ale trupului său, din carnea și oasele sale.

\VS{31}“Pentru aceasta, omul va lăsa pe tatăl său și pe mama sa și se va uni cu soția sa. Atunci cei doi vor deveni un singur trup.”

\VS{32}Acest mister este mare, dar eu vorbesc despre Cristos și despre Adunare.

\VS{33}Cu toate acestea, fiecare dintre voi trebuie să-și iubească și el soția ca pe sine însuși; iar soția să vadă că își respectă soțul.



\par }\Chap{6}{\PP \VerseOne{1}Copii, ascultați de părinții voștri în Domnul, căci așa este drept.

\VS{2}“Cinstiți pe tatăl vostru și pe mama voastră”, care este prima poruncă cu o promisiune:

\VS{3}“ca să vă fie bine și să trăiți mult pe pământ”.


\par }{\PP \VS{4}Voi, taților, nu vă mâniați copiii, ci creșteți-i în disciplina și învățătura Domnului.


\par }{\PP \VS{5}Slujitori, fiți supuși celor care, după trup, sunt stăpânii voștri, cu frică și cu cutremur, în curăția inimii voastre, ca față de Hristos,

\VS{6}nu în felul de a sluji numai când ochii sunt ațintiți asupra voastră, ca niște plăceri de oameni, ci ca niște slujitori ai lui Hristos, făcând din inimă voia lui Dumnezeu,

\VS{7}cu bunăvoință, slujind ca Domnului și nu ca oamenilor,

\VS{8}știind că orice lucru bun pe care îl face fiecare, va primi din nou același bine de la Domnul, fie că este legat, fie că este liber.


\par }{\PP \VS{9}Voi, stăpânilor, faceți-le și voi la fel și lăsați-vă amenințați, știind că Cel ce este și stăpânul lor și al vostru este în ceruri și că la El nu este nici un fel de părtinire.


\par }{\PP \VS{10}În sfârșit, întăriți-vă în Domnul și în tăria puterii Lui.

\VS{11}Îmbrăcați-vă cu toată armura lui Dumnezeu, ca să puteți sta împotriva vicleniei diavolului.

\VS{12}Căci lupta noastră nu este împotriva cărnii și a sângelui, ci împotriva principatelor, împotriva puterilor, împotriva stăpânitorilor lumii, a întunericului acestui veac și împotriva forțelor spirituale ale răutății din locurile cerești.

\VS{13}De aceea, îmbrăcați-vă cu toată armura lui Dumnezeu, ca să puteți rezista în ziua cea rea și, după ce ați făcut totul, să rămâneți în picioare.

\VS{14}Rămâneți deci în picioare, având centura utilă a adevărului încheiată la brâu, îmbrăcând platoșa neprihănirii,

\VS{15}și încălțându-vă picioarele cu pregătirea Bunei Vestiri a păcii,

\VS{16}și, mai presus de toate, luând scutul credinței, cu care veți putea stinge toate săgețile aprinse ale celui rău.

\VS{17}și luați coiful mântuirii și sabia Duhului, care este Cuvântul lui Dumnezeu;

\VS{18}cu toată rugăciunea și cererile, rugându-vă în orice moment în Duhul și veghind în acest scop cu toată stăruința și cererile pentru toți sfinții.

\VS{19}Rugați-vă pentru mine, ca să mi se dea posibilitatea de a-mi deschide gura, ca să fac cunoscut cu îndrăzneală misterul Bunei Vestiri,

\VS{20}pentru care sunt ambasador în lanțuri; ca în el să vorbesc cu îndrăzneală, așa cum trebuie să vorbesc.


\par }{\PP \VS{21}Dar, ca să știți și voi cum stau lucrurile cu mine și cum mă descurc, Tihic, fratele preaiubit și slujitorul credincios în Domnul, vă va face cunoscute toate lucrurile.

\VS{22}Tocmai de aceea vi l-am trimis la voi, ca să cunoașteți starea noastră și ca să vă mângâie inimile.


\par }{\PP \VS{23}Pace fraților și dragoste cu credință, de la Dumnezeu Tatăl și de la Domnul Isus Hristos.

\VS{24}Harul să fie cu toți cei ce iubesc pe Domnul nostru Isus Hristos cu dragoste incoruptibilă. Amin.


\par }